%% ===== 2 Background =====
\section{Background}
\label{sec:background}

% [27/10] Background chỉ định nghĩa; nguồn gốc hai bảng phân loại và quy tắc mã ở Method 4.5, độ nhạy của định nghĩa legacy
% ở Results 5.3. Ví dụ kiểm với bài gốc (FreshBrew, Google, RISE, GUIMigrator); legacy theo công nghệ nguồn (comment 4 của
% Hà Cao); thêm câu migration là từ bao trùm, translation là XL và DIA.
\paragraph{Kinds of migration.} \new{We distinguish five kinds of migration:} between programming languages (XL, e.g.\ C to Rust \new{or} COBOL to Java); between versions of a language or platform (VER, e.g.\ Java~8 to~17); between libraries or APIs, including dependency upgrades with breaking changes (LIB, e.g.\ JUnit~3 to~4); between dialects (DIA, e.g.\ PostgreSQL to MySQL); and between frameworks or application platforms (FW, e.g.\ \new{Android XML views to SwiftUI}). \new{We use migration as the umbrella term and call the XL and DIA kinds translation, because they rewrite code in another language or dialect. A study may involve several kinds. Independently of the kind, a \emph{legacy} flag marks a source written for a language or platform that is no longer in mainstream use for new development, such as COBOL, PL/I, mainframe systems or ASP.NET Web Forms. The flag follows the source technology, not the authors' description of their code.} \cmt{Claude}{[R07] This example comes from GUIMigrator, which targets both Jetpack Compose and SwiftUI. The SwiftUI target also changes language and platform (to Swift on iOS), and the extraction sheet itself notes that FW/XL ``could be considered''; yet GUIMigrator is coded FW only, while IP-preserving is coded FW/XL. Decide one rule, and use the Compose target (Android XML views to Jetpack Compose) as the FW example.}

\begin{added}
% [27/10] Danh sách có dấu phẩy bên trong nên tách bằng chấm phẩy; "run" thành "execute" như RQ4.
\paragraph{Behavioral equivalence.} Two programs are behaviorally equivalent if they produce the same observable behavior for every input: the same outputs; the same effects on files, databases or other state; and the same termination or failure. Checking this exactly is undecidable in general, so every practical check approximates it. Tests and differential runs sample the input space; bounded model checking and symbolic execution cover it up to a bound; formal proofs cover it completely but only for a model of the program; and compilation, similarity metrics and human or LLM judgment do not execute the program. Migration also relaxes this definition in ways that an evaluation should make explicit. A translation may map types whose ranges differ, such as \new{32-bit integers in C into Python's unbounded integers}; a library migration may keep the functional result but change logging, performance or error messages; and a framework migration may be judged on what a user sees rather than on internal state. ``Equivalent'' therefore always means equivalent with respect to some oracle and some notion of observable behavior.
\end{added}

% [27/10] Oracle là "the check a study uses" vì O-a, O-b không trực tiếp xét hành vi. Bảng 1: sửa O-b, O-d, O-e, O-f, O-h
% cho bớt tuyệt đối và rõ nghĩa; cột "What it checks" rộng hơn (0,24/0,32/0,28). Trục nguồn của kết quả mong đợi theo
% comment 5 của Hà Cao; BatFix vẫn là O-c vì test của TransCoder có trước nghiên cứu.
\paragraph{Oracles.} \new{An oracle is the check a study uses to decide} whether the migrated program behaves like the original. We use nine categories (Table~\ref{tab:oracles}).

\begin{table*}[t]
  \centering
  \caption{Oracle categories used throughout the paper.}
  \label{tab:oracles}
  \small
  \begin{tabular}{@{}l P{0.24\linewidth} P{0.32\linewidth} P{0.28\linewidth}@{}}
    \toprule
    Code & Oracle & What it checks & Known weakness \\
    \midrule
    O-a & Surface similarity (BLEU, CodeBLEU, exact match, AST match) & Closeness to a reference translation & Equivalent code can look different, and similar code can behave differently \\
    O-b & Compilation, type checking \new{or another well-formedness check} & The output is well-formed in the target & \new{Does not check behavior} \\
    O-c & Existing tests & \new{Behavior on a test suite that existed before the study, built for a benchmark or by the project's developers} & Only as strong as the tests; benchmarks often have few \\
    O-d & \new{Translated tests} & \new{Behavior on the original program's own tests after they are translated with the code; in test migration, on the migrated tests themselves} & Test and code may share the same translation error\new{; a migrated test that passes may no longer check what the original did} \\
    O-e & LLM-generated tests & \new{Behavior on new tests whose expected outputs come from an LLM or a specification} & \new{Expected outputs can be wrong or share a source with the output under test} \\
    O-f & \new{Differential testing} & \new{Same outputs as the original program, run by the study on the same inputs (from fuzzing, existing tests or an LLM)} & Limited by input generation; needs the original to run \\
    O-g & \new{Formal verification} & Equivalence proven or bounded-checked \new{against the original or a specification} & Costly; bounds and models limit coverage \\
    O-h & \new{LLM judge} & An LLM's verdict on equivalence & \new{The judge's own accuracy is rarely validated} \\
    O-i & Human evaluation & A person's judgment & Subjective, small samples, rarely agreement-measured \\
    \bottomrule
  \end{tabular}
\end{table*}

% [27/10] "source" thành "origin" để khỏi nhầm với chương trình nguồn, thống nhất với "ground-truth origin" và bảng tab:origin ở Results 5.3.
% O-d "translated with the code" thành "translated for the study" cho đối xứng với O-c (quy tắc ở 4.5).
\new{The four test-based categories (O-c to O-f) differ in the origin of their ground truth, the expected outputs: a test suite that existed before the study (O-c); the original program's tests, translated for the study (O-d); an LLM or a specification (O-e); or a run of the original program during the study (O-f). Tests that an LLM writes therefore fall under O-f when their expected outputs come from running the original, because the LLM then supplies only the inputs.} \cmt{Claude}{[R08, major] The label ``LLM-generated tests'' (O-e) names who writes the tests, but the category is defined by where the expected outputs come from (an LLM or a specification): LLM-written tests whose expected outputs come from running the original are O-f (previous sentence), and by Table 1 a test written from a specification without any LLM would be O-e. Readers of the abstract will conclude that LLM-generated tests are rare (4\%), although 3 more papers run LLM-generated inputs under O-f. Consider a label based on the origin of the expected outputs, e.g.\ ``Tests with LLM- or specification-derived expected outputs'', or say this explicitly wherever the 4\% is reported.}

The categories are not exclusive. Many studies combine several, and some use one inside the technique (for example to filter training data or to drive a repair loop) and another to report results. \new{Section~\ref{sec:extraction} states which of these uses we code.}
